\section*{Abstract}
		Dieses Dokument bietet eine umfassende Analyse der fundamentalen Beziehung zwischen dem geometrischen Parameter $\xipar = \frac{4}{3} \times 10^{-4}$ der T0-Theorie und der Euler'schen Zahl $e = 2.71828\ldots$ Die T0-Theorie basiert auf tiefen geometrischen Prinzipien aus tetraedrischer Packung und postuliert eine fraktale Raumzeit mit lokaler Dimension $D_f^{\,\text{Raum}} = 3 - \xi \approx 2{,}9999$ und einer effektiven Dimension $D_f^{\,\text{eff}} \approx 2{,}973$ aus dem kumulativen Skalenfluss. Wir zeigen detailliert, wie Potenzen und Exponentialformen von $\xipar$ die Hierarchie der Teilchenmassen (als Verhältnisse), Zeitskalen und fundamentalen Konstanten aus ersten Prinzipien beschreiben. Besonderes Augenmerk liegt auf der mathematischen Konsistenz und den experimentell überprüfbaren Vorhersagen der Theorie.

	
	
	\section{Einleitung: Die geometrische Basis der T0-Theorie}
	
	\subsection{Historische und konzeptionelle Grundlagen}
	
	Die T0-Theorie entstand aus der Beobachtung, dass fundamentale physikalische Konstanten und Massenverhältnisse nicht zufällig verteilt sind, sondern tiefen mathematischen Beziehungen folgen. Im Gegensatz zu vielen anderen Ansätzen postuliert T0 keine neuen Teilchen oder zusätzlichen Dimensionen, sondern eine fundamentale geometrische Struktur der Raumzeit selbst.
	
	\begin{erkenntnis}
		\textbf{Das zentrale Paradigma der T0-Theorie:}
		
		Die Physik auf fundamentaler Ebene ist nicht durch zufällige Parameter charakterisiert, sondern durch eine zugrundeliegende geometrische Struktur, die durch den Parameter $\xi$ quantifiziert wird. Die Euler'sche Zahl $e$ dient als der natürliche Operator, der diese geometrische Struktur in dynamische Prozesse übersetzt.
	\end{erkenntnis}
	
	\subsection{Die tetraedrische Herkunft von $\xi$}
	
	\begin{beziehung}
		\textbf{Geometrische Ableitung von $\xi = \frac{4}{3} \times 10^{-4}$:}
		
		Die fundamentale Konstante $\xi$ leitet sich aus der Geometrie regelmäßiger Tetraeder ab. Für einen Tetraeder mit Kantenlänge $a$:
		
		\begin{align}
			V_{\text{tetra}} &= \frac{\sqrt{2}}{12}a^3 \\
			R_{\text{umkugel}} &= \frac{\sqrt{6}}{4}a \\
			V_{\text{sphäre}} &= \frac{4}{3}\pi R_{\text{umkugel}}^3 = \frac{\pi\sqrt{6}}{8}a^3 \\
			\frac{V_{\text{tetra}}}{V_{\text{sphäre}}} &= \frac{\sqrt{2}/12}{\pi\sqrt{6}/8} = \frac{2\sqrt{3}}{9\pi} \approx 0.1225
		\end{align}
		
		Durch Skalierung und Normierung ergibt sich:
		\begin{equation}
			\xipar = \frac{4}{3} \times 10^{-4} = \left(\frac{V_{\text{tetra}}}{V_{\text{sphäre}}}\right) \times \text{Skalierungsfaktor}
		\end{equation}
		
		\begin{center}
			\begin{tikzpicture}[scale=1.4]
				% Regelmäßiges Tetraeder
				\coordinate (A) at (0,0);
				\coordinate (B) at (2,0);
				\coordinate (C) at (1,1.732);
				\coordinate (D) at (1,0.577);
				
				\draw[t0blue, thick] (A) -- (B) -- (C) -- cycle;
				\draw[t0blue, thick] (A) -- (D);
				\draw[t0blue, thick] (B) -- (D);
				\draw[t0blue, thick] (C) -- (D);
				
				% Umschriebene Kugel
				\draw[t0red, dashed] (1,0.577) circle (1.155);
				
				\node at (0,0) [below left] {A};
				\node at (2,0) [below right] {B};
				\node at (1,1.732) [above] {C};
				\node at (1,0.577) [below] {D (Schwerpunkt)};
				
				\node at (3.2,0.866) [t0blue, align=left] {Tetraeder: $V = \frac{\sqrt{2}}{12}a^3$};
				\node at (3.2,0.5) [t0red, align=left] {Umkugel: $V = \frac{\pi\sqrt{6}}{8}a^3$};
			\end{tikzpicture}
		\end{center}
	\end{beziehung}
	
	\subsection{Die fraktale Raumzeit-Dimension}
	
	\begin{abhandlung}
		\textbf{Die fraktale Natur der Raumzeit}
		
		Eine der radikalsten Aussagen der T0-Theorie ist, dass die Raumzeit auf fundamentaler Ebene fraktale Eigenschaften besitzt. Die effektive Dimension hängt von der Energieskala ab:
		
		\begin{equation}
			D_f(E) = 4 - 2\xipar \cdot \ln\left(\frac{E_P}{E}\right)
		\end{equation}
		
				Diese Formel ergibt bei $E = E_P$ genau $4$ und fällt zu kleinen Energien hin nur schwach ab ($3{,}988$ bei $1$~GeV, $3{,}983$ bei $1$~eV); ein Wert von $2{,}973$ verlangte $\ln(E_P/E) \approx 3850$ und folgt aus ihr nicht. Die effektive fraktale Dimension stammt aus dem kumulativen Skalenfluss (A040):
\begin{equation}
  \resizebox{\textwidth}{!}{$\displaystyle D_f^{\,\text{eff}} \approx 2{,}973 \quad \text{(effektive fraktale Dimension über kumulativen Skalenfluss)}$}
\end{equation}
		
		\textbf{Physikalische Interpretation:}
		\begin{itemize}
			\item Bei kleinen Abständen/hohen Energien wird die fraktale Struktur der Raumzeit sichtbar
			\item Die effektive Dimension $D_f^{\,\text{eff}} \approx 2{,}973$ ist kein Zufall, sondern folgt aus der kumulativen Wirkung der fraktalen Struktur über den Skalenfluss
			\item Dies erklärt das Renormierungsverhalten der Quantenfeldtheorien
		\end{itemize}
		
		Die Näherungsformel
		\begin{equation}
			2 + \frac{\ln(1/\xipar)}{\ln(E_P/E_0)} = 2 + \frac{\ln 7500}{\ln(1{,}221\times10^{19})} = 2{,}203
		\end{equation}
		mit $E_P = 1.221 \times 10^{19}$ GeV (Planck-Energie) und $E_0 = 1$ GeV (Referenzenergie) liefert den Wert $2{,}973$ ebenfalls nicht.
	\end{abhandlung}
	
	\section{Die Euler'sche Zahl als dynamischer Operator}
	
	\subsection{Mathematische Grundlagen von $e$}
	
	\begin{beziehung}
		\textbf{Die einzigartigen Eigenschaften von $e$:}
		
		Die Euler'sche Zahl ist durch mehrere äquivalente Definitionen charakterisiert:
		
		\begin{align}
			e &= \lim_{n \to \infty} \left(1 + \frac{1}{n}\right)^n \\
			e &= \sum_{n=0}^{\infty} \frac{1}{n!} \\
			\frac{d}{dx}e^x &= e^x \\
			\int e^x dx &= e^x + C
		\end{align}
		
		In der T0-Theorie erhält $e$ eine besondere Bedeutung als der natürliche Übersetzer zwischen diskreter geometrischer Struktur und kontinuierlicher dynamischer Entwicklung.
	\end{beziehung}
	
	\subsection{Zeit-Masse-Dualität als fundamentales Prinzip}
	
	\begin{erkenntnis}
		\textbf{Die Zeit-Masse-Dualität: $T \cdot m = 1$}
		
		In natürlichen Einheiten ($\hbar = c = 1$) gilt die fundamentale Beziehung:
		\begin{equation}
			\boxed{T \cdot m = 1}
		\end{equation}
		
		Dies bedeutet:
		\begin{itemize}
			\item Jedes Teilchen hat eine charakteristische Zeitskala $T = 1/m$
			\item Schwere Teilchen leben typischerweise kürzer
			\item Leichte Teilchen haben längere charakteristische Zeitskalen
			\item Die $\xi$-Modulation führt zu Korrekturen der Ordnung $\xipar$
		\end{itemize}
		
		\textbf{Beispiele:}
		\begin{align}
			\text{Elektron: } & T_e = \hbar/(m_e c^2) \approx 1.29 \times 10^{-21}\, \text{s} \\
			\text{Myon: } & T_\mu = \hbar/(m_\mu c^2) \approx 6.23 \times 10^{-24}\, \text{s} \\
			\text{Tauon: } & T_\tau = \hbar/(m_\tau c^2) \approx 3.70 \times 10^{-25}\, \text{s}
		\end{align}
		
		Diese Zeitskalen sind die Compton-Zeiten der Leptonen. Sie sind nicht die Lebensdauern: Das Myon lebt $2.2 \times 10^{-6}$~s, das Tauon $2.9 \times 10^{-13}$~s, das Elektron ist stabil.
		
	\end{erkenntnis}
	
	\section{Detaillierte Analyse der Leptonenmassen}
	
	\subsection{Die Massenhierarchie als Verhältnisse}
	
	\begin{beziehung}
		\textbf{Die Leptonenhierarchie in Verhältnissen:}
		
		Tragend sind nicht die absoluten Massen, sondern die Verhältnisse zwischen den Generationen. Sie folgen als Potenzen von $\xipar$ mit rationalen Vorfaktoren (Leptonleiter, Dok.~352):
		\begin{align}
			\frac{m_\mu}{m_e} &= \frac{12}{5}\,\xipar^{-1/2} = \frac{12}{5}\sqrt{7500} = 207.85 \\
			\frac{m_\tau}{m_\mu} &= \frac{125}{144}\,\xipar^{-1/3} = 16.99 \\
			\frac{m_\tau}{m_e} &= \frac{25}{12}\,\xipar^{-5/6} = 3531.6
		\end{align}
		
		Vergleich mit der Messung:
		\begin{align}
			\text{Experimentell: } \frac{m_\mu}{m_e} &= 206.77 \quad (+0.52\,\%) \\
			\frac{m_\tau}{m_\mu} &= 16.817 \quad (+1.04\,\%) \\
			\frac{m_\tau}{m_e} &= 3477.2 \quad (+1.56\,\%)
		\end{align}
		
		Die berechneten Werte sind bare-Verhältnisse; die verbleibenden Abweichungen sind die generationsabhängigen Reste zwischen bare- und Pol-Masse (Dok.~352). Eine absolute Masse ist keine eigene Vorhersage: Sie ergibt sich aus dem Verhältnis und einem einzigen Anker (z.\,B. $m_e$).
		
	\end{beziehung}
	
	\subsection{Logarithmische Symmetrie und ihre Konsequenzen}
	
	\begin{abhandlung}
		\textbf{Die tiefere Bedeutung der logarithmischen Symmetrie:}
		
		Die Beziehung $\ln(m_\mu) = \frac{\ln(m_e) + \ln(m_\tau)}{2}$, also $m_\mu = \sqrt{m_e m_\tau}$, gilt für die gemessenen Massen nicht (sie ergibt $30.1$~MeV). Sie gilt erst mit dem Zusatzfaktor der Leptonleiter:
		\begin{equation}
			m_\mu = \frac{24\sqrt{3}}{25}\,\xipar^{-1/12}\sqrt{m_e \cdot m_\tau} = 105.4\ \text{MeV}
		\end{equation}

		Der Faktor $\frac{24\sqrt3}{25}\xi^{-1/12} = 3.497$ folgt aus $m_\mu/m_e = \frac{12}{5}\xi^{-1/2}$ und $m_\tau/m_\mu = \frac{125}{144}\xi^{-1/3}$; $105.4$~MeV ist der bare-Wert, $-0.26\,\%$ gegen die Messung (Dok.~352).
		
		Dies ist keine zufällige Koinzidenz, sondern weist auf eine zugrundeliegende algebraische Struktur hin. In der Gruppen-theoretischen Interpretation entsprechen die Leptonen verschiedenen Darstellungen einer zugrundeliegenden Symmetrie.
		
		\textbf{Mögliche Interpretationen:}
		\begin{itemize}
			\item Die Leptonen entsprechen verschiedenen Energielevel in einem geometrischen Potential
			\item Es gibt eine diskrete Skalierungssymmetrie mit den Stufen $\xipar^{-1/2}$ und $\xipar^{-1/3}$
			\item Die Exponenten der Leptonleiter könnten mit topologischen Ladungen zusammenhängen
		\end{itemize}
		
		Die Konsistenz über drei Generationen hinweg ist bemerkenswert und spricht gegen Zufall.
	\end{abhandlung}
	
	\section{Fraktale Raumzeit und Quantenfeldtheorie}
	
	\subsection{Das Renormierungsproblem und seine Lösung}
	
	\begin{anwendung}
		\textbf{Die T0-Lösung der UV-Divergenzen:}
		
		In konventioneller Quantenfeldtheorie treten Divergenzen auf wie:
		\begin{equation}
			\int_0^\infty \frac{d^4k}{k^2 - m^2} \to \infty
		\end{equation}
		
		Die fraktale Raumzeit mit $D_f^{\,\text{eff}} \approx 2{,}973$ führt zu einem natürlichen Cutoff:
		\begin{equation}
			\boxed{\Lambda_{\text{T0}} = \frac{E_P}{\xipar} \approx 9.2 \times 10^{22}\, \text{GeV}}
		\end{equation}

		
		Propagator-Modifikation:
		\begin{equation}
			G(k) = \frac{1}{k^2 - m^2} \cdot e^{-\xipar \cdot k/E_P}
		\end{equation}
		
		\textbf{Wirkung auf Feynman-Diagramme:}
		\begin{itemize}
			\item Schleifenintegrale werden natürlich regularisiert
			\item Keine willkürlichen Cutoffs notwendig
			\item Die Regularisierung ist lorentzinvariant
			\item Skalenfluss ($\mathcal{R}$-induziert) wird modifiziert
		\end{itemize}
		
		\begin{equation}
			\int_0^\infty d^4k\, G(k) \cdot e^{-\xipar \cdot k/E_P} < \infty
		\end{equation}
	\end{anwendung}
	
	\subsection{Modifizierte Skalenflussequationen}
	
	\begin{beziehung}
		\textbf{Skalenfluss in fraktaler Raumzeit ($\mathcal{R}$-induziert):}
		
		Die beta-Funktion für die Kopplungskonstante $\alpha$ wird modifiziert:
		\begin{equation}
			\frac{d\alpha}{d\ln\mu} = \beta_0 \alpha^2 \cdot \left(1 + \xipar \cdot \ln\frac{\mu}{E_0}\right)
		\end{equation}
		
		Für die Feinstrukturkonstante:
		\begin{equation}
			\alpha^{-1}(\mu) = \alpha^{-1}(m_e) - \beta_0 L - \beta_0 \xipar \left(\frac{L^2}{2} + L \ln\frac{m_e}{E_0}\right), \qquad L = \ln\frac{\mu}{m_e}
		\end{equation}
		Die Form folgt aus $d\alpha^{-1}/d\ln\mu = -\beta_0(1+\xi\ln(\mu/E_0))$; der Kreuzterm $L\ln(m_e/E_0)$ tritt auf, weil der Bezugspunkt $E_0 \neq m_e$ ist ($\ln(m_e/E_0) = -2{,}67$).
		
		\textbf{Konsequenzen:}
		\begin{itemize}
			\item Leichte Modifikation der laufenden Kopplungen
			\item Vorhersage von kleinen Abweichungen bei hohen Energien
			\item Testbar an LHC-Daten
		\end{itemize}
	\end{beziehung}
	
	\section{Kosmologische Anwendungen und Vorhersagen}
	
	\subsection{Urknall und CMB-Temperatur}
	
	\begin{anwendung}
		\textbf{Ansatz für die CMB-Temperatur:}

		Der Ansatz für die heutige Temperatur der kosmischen Hintergrundstrahlung lautet:
		\begin{equation}
			T_{\text{CMB}} = T_P \cdot e^{-\xipar \cdot N}
		\end{equation}
		
		Mit:
		\begin{itemize}
			\item $T_P = 1.416 \times 10^{32}$ K (Planck-Temperatur)
			\item $N = 114$ (Anzahl der $\xi$-Skalierungen)
			\item $\xipar \cdot N = 1.333 \times 10^{-4} \times 114 = 0.0152$
		\end{itemize}
		
		Berechnung:
		\begin{align}
			T_{\text{CMB}} &= 1.416 \times 10^{32} \cdot e^{-0.0152} \\
			&= 1.416 \times 10^{32} \cdot 0.9849 \\
			&= 1.395 \times 10^{32}\, \text{K}
		\end{align}
		
		Der Wert liegt $32$ Größenordnungen über den gemessenen $2.725$~K. Eine Herleitung der CMB-Temperatur aus $T_P\,e^{-\xi N}$ gilt daher nicht: Mit $N = 114$ lässt sich $T_{\text{CMB}}$ nicht erreichen; $2.725$~K verlangte $\xi N = 73.0$, also $N \approx 5.5\times10^{5}$.
	\end{anwendung}
	
	\subsection{Dunkle Energie und kosmologische Konstante}
	
	\begin{erkenntnis}
		\textbf{Das dunkle Energie-Problem gelöst?}
		
		Die Vakuumenergiedichte in T0:
		\begin{equation}
			\rho_{\Lambda} = \frac{E_P^4}{(2\pi)^3} \cdot \xipar^2
		\end{equation}
		
		Numerisch:
		\begin{align}
			E_P^4 &= (1.221 \times 10^{19}\, \text{GeV})^4 = 2.23 \times 10^{76}\, \text{GeV}^4 \\
			\xipar^2 &= (1.333 \times 10^{-4})^2 = 1.777 \times 10^{-8} \\
			\rho_{\Lambda} &\approx 3.96 \times 10^{68} \cdot 1.777 \times 10^{-8} = 7.04 \times 10^{60}\, \text{GeV}^4
		\end{align}
		
		Umrechnung in beobachtbare Einheiten:
		\begin{equation}
			\rho_{\Lambda} \approx 10^{-123} E_P^4
		\end{equation}
		
		\textbf{Genau in der richtigen Größenordnung für dunkle Energie!}
		
		Die T0-Theorie erklärt natürlicherweise, warum die Vakuumenergiedichte so unglaublich klein ist im Vergleich zur Planck-Skala.
	\end{erkenntnis}
	
	\section{Experimentelle Tests und Vorhersagen}
	
	\subsection{Präzisionstests in der Teilchenphysik}
	
	\begin{anwendung}
		\textbf{Spezifische, testbare Vorhersagen:}
		
		\begin{enumerate}
			\item \textbf{Leptonen-Massenverhältnis:}
			\begin{equation}
				\frac{m_\mu}{m_e} = 206.768282 \cdot (1 + \alpha \xipar + \beta \xipar^2 + \cdots)
			\end{equation}
			Abweichungen bei 0.01\%-Präzision messbar
			
			\item \textbf{Neutrino-Oszillationen:}
			\begin{equation}
				P(\nu_\alpha \to \nu_\beta) = P_{\text{SM}} \cdot (1 + \gamma \xipar \cdot L/E)
			\end{equation}
			Modifikation der Oszillationswahrscheinlichkeit
			
			\item \textbf{Myon-Zerfall:}
			\begin{equation}
				\Gamma(\mu \to e\nu_e\nu_\mu) = \Gamma_{\text{SM}} \cdot e^{-\xipar \cdot m_\mu/E_P}
			\end{equation}
			Kleine Korrekturen zur Zerfallsrate
			
			\item \textbf{Anomales magnetisches Moment:}
			\begin{equation}
				a_e = a_e^{\text{SM}} \cdot (1 + \delta \xipar)
			\end{equation}
			Erklärung der möglichen Anomalien
		\end{enumerate}
	\end{anwendung}
	
	\subsection{Kosmologische Tests}
	
	\begin{anwendung}
		\textbf{Tests mit kosmologischen Daten:}
		
		\begin{itemize}
			\item \textbf{CMB-Spektrum:} Vorhersage spezifischer Modifikationen des CMB-Leistungsspektrums aufgrund der fraktalen Raumzeit
			
			\item \textbf{Strukturbildung:} Modifiziertes Skalierungsverhalten der Materieverteilung
			
			\item \textbf{Primordiale Nucleosynthese:} Leichte Modifikationen der Elementhäufigkeiten aufgrund geänderter Expansionsrate im frühen Universum
			
			\item \textbf{Gravitationswellen:} Vorhersage einer skalaren Komponente in primordialen Gravitationswellen
		\end{itemize}
		
		\begin{equation}
			h_{\mu\nu} = h_{\mu\nu}^{\text{tensor}} + \xipar \cdot h^{\text{skalar}}
		\end{equation}
	\end{anwendung}
	
	\section{Mathematische Vertiefung}
	
	\subsection{Die $\pi$-$e$-$\xi$ Trinität}
	
	\begin{beziehung}
		\textbf{Die fundamentale Dreiheit:}
		
		Die drei mathematischen Konstanten $\pi$, $e$ und $\xi$ spielen komplementäre Rollen:
		
		\begin{align}
			\pi &: \text{Geometrie und Topologie} \\
			e &: \text{Wachstum und Dynamik} \\
			\xi &: \text{Kopplung und Skalierung}
		\end{align}
		
		Ihre Kombination erscheint in fundamentalen Beziehungen:
		
		\begin{equation}
			e^{i\pi} + 1 = 0 \quad \text{(klassische Euler-Identität)}
		\end{equation}
		
		\begin{equation}
			e^{i\xipar\pi} + 1 \approx \delta(\xipar) \quad \text{(T0-Erweiterung)}
		\end{equation}
		
		\begin{equation}
			\frac{m_\mu}{m_e} = \frac{12}{5}\,\xipar^{-1/2} \quad \text{(Massenhierarchie als Verhältnis)}
		\end{equation}
		
		\begin{center}
			\begin{tikzpicture}[scale=2.2]
				\draw[thick, t0blue] (0,0) circle (1);
				\node at (90:1.3) [t0blue, align=center] { $\pi$ \\ \small Geometrie \\ \small Symmetrie};
				
				\node at (210:1.3) [t0green, align=center] { $e$ \\ \small Dynamik \\ \small Wachstum};
				
				\node at (330:1.3) [t0orange, align=center] { $\xi$ \\ \small Kopplung \\ \small Quantisierung};
				
				\draw[->, thick, t0blue] (90:0.8) -- (210:0.8);
				\draw[->, thick, t0green] (210:0.8) -- (330:0.8);
				\draw[->, thick, t0orange] (330:0.8) -- (90:0.8);
				
				\node at (0,0) {$e^{i\xi\pi}$};
			\end{tikzpicture}
		\end{center}
	\end{beziehung}
	
	\subsection{Gruppentheoretische Interpretation}
	
	\begin{abhandlung}
		\textbf{Mögliche gruppentheoretische Basis:}
		
		Die Exponenten der Leptonleiter, $-\tfrac12$ für $m_\mu/m_e$ und $-\tfrac13$ für $m_\tau/m_\mu$, sind Vielfache von $\tfrac16$ (zusammen $-\tfrac56$). Das legt nahe, dass die Leptonen-Generationen mit Darstellungen einer diskreten Gruppe zusammenhängen, deren Ordnung die Stufen $\tfrac16$ festlegt; die Herleitung der Vorfaktoren steht in Dok.~352.
		
		
		\textbf{Mögliche Interpretation:}
		Die Leptonen-Generationen entsprechen verschiedenen Ladungen unter einer diskreten Eichsymmetrie, die aus der zugrundeliegenden geometrischen Struktur emergiert.
	\end{abhandlung}

	\section{Experimentelle Konsequenzen}
	
	\subsection{Präzisionsvorhersagen}
	
	\begin{anwendung}
		\textbf{Testbare Vorhersagen:}
		
		\begin{enumerate}
			\item \textbf{Leptonen-Verhältnis:}
			\begin{equation}
				\frac{m_\mu}{m_e} = 206.768282 \cdot (1 + \alpha \xi + \beta \xi^2 + \cdots)
			\end{equation}
			
			\item \textbf{Myon-Zerfall:}
			\begin{equation}
				\Gamma(\mu \to e\nu_e\nu_\mu) = \Gamma_{\text{SM}} \cdot e^{-\xi \cdot m_\mu/E_P}
			\end{equation}
			
			\item \textbf{Anomales magnetisches Moment:}
			\begin{equation}
				a_e = a_e^{\text{SM}} \cdot (1 + \delta \xi)
			\end{equation}
			
			\item \textbf{Neutrino-Oszillationen:}
			\begin{equation}
				P(\nu_\alpha \to \nu_\beta) = P_{\text{SM}} \cdot (1 + \gamma \xi \cdot L/E)
			\end{equation}
		\end{enumerate}
	\end{anwendung}
	